\documentclass[a4paper,11pt]{article}
\usepackage[utf8]{inputenc}
\usepackage[T1]{fontenc}
\usepackage[polish]{babel}
\usepackage[margin=2cm]{geometry}
\usepackage{enumitem}
\usepackage{titlesec}
\usepackage{parskip}
\setlength{\headheight}{14pt}
\usepackage{fancyhdr}
\usepackage{booktabs}
\usepackage{array}

% Nagłówek i stopka
\pagestyle{fancy}
\fancyhf{}
\rhead{\textit{Projekt skryptu — meteo}}
\lhead{\textit{[Piotr Siebierański]}}
\cfoot{\thepage}

% Formatowanie sekcji
\titleformat{\section}{\large\bfseries}{}{0em}{}[\titlerule]
\titlespacing{\section}{0pt}{12pt}{6pt}

\begin{document}

% ───────────────────────────── NAGŁÓWEK ─────────────────────────────
\begin{center}
    {\LARGE\bfseries Projekt skryptu: \texttt{meteo}}\\[6pt]
    {\large [Piotr Siebierański]}\\[2pt]
    {\small Systemy operacyjne — projekt dużego skryptu}
\end{center}

\vspace{4pt}
\hrule
\vspace{12pt}

% ───────────────────────────── OPIS FUNKCJONALNOŚCI ─────────────────
\section{Opis funkcjonalności}

Skrypt \texttt{meteo} pobiera prognozę pogody dla podanej lokalizacji
i prezentuje ją użytkownikowi w postaci wykresu graficznego generowanego
przez program \texttt{gnuplot}. Skrypt umożliwia:

\begin{itemize}[noitemsep]
    \item pobieranie prognozy na 1--3 dni do przodu z serwisu \texttt{wttr.in},
    \item wybór prezentowanych danych: temperatura, opady, prędkość wiatru,
    \item generowanie wykresu PNG otwieranego automatycznie przez \texttt{xdg-open},
    \item opcjonalny zapis wykresu do wskazanego pliku bez otwierania przeglądarki,
    \item konfigurację domyślnych parametrów przez plik \texttt{.rc},
    \item nadpisanie dowolnego parametru z linii poleceń.
\end{itemize}

% ───────────────────────────── OPIS DZIAŁANIA ───────────────────────
\section{Opis działania}


\begin{enumerate}[noitemsep]
    \item Skrypt wczytuje plik konfiguracyjny (domyślnie \texttt{\textasciitilde/.meteorc},
          możliwy do nadpisania flagą \texttt{-c}).
    \item Parametry z linii poleceń nadpisują wartości z konfiguracji.
    \item Skrypt tworzy katalog tymczasowy \texttt{/tmp/meteo\_\$\$} (gdzie \texttt{\$\$}
          to PID procesu); po zakończeniu katalog jest usuwany przez \texttt{trap EXIT}.
    \item Dane pogodowe są pobierane z API \texttt{wttr.in} w formacie JSON
          i parsowane narzędziem \texttt{jq}.
    \item Na podstawie pobranych danych generowany jest plik danych dla
          \texttt{gnuplot}, a następnie wykres PNG.
    \item Wykres jest otwierany przez \texttt{xdg-open} lub zapisywany do pliku,
          jeśli podano flagę \texttt{-o}.
\end{enumerate}

\subsection*{Interfejs użytkownika}

Skrypt działa wyłącznie z poziomu linii poleceń. Składnia wywołania:

\begin{center}
\texttt{meteo [opcje]}
\end{center}

\vspace{4pt}
\begin{tabular}{>{\ttfamily}lp{10cm}}
\toprule
\textrm{Flaga} & Opis \\
\midrule
-l \textit{lokalizacja} & Nazwa miasta lub współrzędne (np. \texttt{Warszawa}, \texttt{51.1,17.0}) \\
-d \textit{dni}         & Liczba dni prognozy (1--3, domyślnie z \texttt{.rc}) \\
-t                      & Pokaż wykres temperatury \\
-p                      & Pokaż wykres opadów \\
-w                      & Pokaż wykres prędkości wiatru \\
-o \textit{plik.png}    & Zapisz wykres do pliku zamiast otwierać \\
-c \textit{plik}        & Użyj podanego pliku konfiguracyjnego \\
-v                      & Wyświetl wersję i autora \\
-h                      & Wyświetl krótką pomoc \\
\bottomrule
\end{tabular}

\subsection*{Plik konfiguracyjny \texttt{.meteorc}}

Plik konfiguracyjny przechowuje domyślne wartości parametrów.
Przykładowa zawartość:

\begin{verbatim}
LOKALIZACJA="Warszawa"
JEDNOSTKI="C"          # C = Celsjusz, F = Fahrenheit
DNI=3
GNUPLOT_PATH=""        # puste = szukaj w PATH
JQ_PATH=""             # puste = szukaj w PATH
\end{verbatim}

Jeśli \texttt{GNUPLOT\_PATH} lub \texttt{JQ\_PATH} są puste, skrypt
wykrywa ścieżki automatycznie przez \texttt{which}. Jeśli narzędzie
nie zostanie znalezione, skrypt kończy działanie z czytelnym komunikatem błędu.

\subsection*{Zależności zewnętrzne}

\begin{itemize}[noitemsep]
    \item \texttt{curl} — pobieranie danych z \texttt{wttr.in},
    \item \texttt{jq} — parsowanie odpowiedzi JSON,
    \item \texttt{gnuplot} — generowanie wykresu PNG,
    \item \texttt{xdg-open} — otwieranie wykresu.
\end{itemize}

\end{document}
